\documentclass[11pt,a4paper]{article}
\usepackage[utf8]{inputenc}
\usepackage[vietnamese]{babel}
\usepackage{amsmath,amssymb}
\usepackage{geometry}
\geometry{a4paper, total={180mm,260mm}, left=15mm, right=15mm, top=18mm, bottom=18mm}

\usepackage{tikz}
\usetikzlibrary{shapes.geometric, arrows.meta, positioning, calc}
\tikzset{
  arr/.style={-{Stealth[scale=0.8]}},
  every node/.append style={align=center}
}
\usepackage{tcolorbox}
\tcbuselibrary{skins,breakable}
\usepackage{fontawesome}
\usepackage{fancyhdr}
\usepackage{enumitem}

\pagestyle{fancy}
\fancyhf{}
\renewcommand{\headrulewidth}{0.5pt}
\renewcommand{\footrulewidth}{0.5pt}
\fancyhead[L]{\small\bfseries\color{violet!85!black}\faBook\ SỔ TAY KIẾN THỨC TOÁN 2}
\fancyhead[R]{\small\bfseries\color{violet!85!black}Lớp Toán Thầy Lĩnh -- loptoanthaylinh.vn}
\fancyfoot[C]{\small\color{gray}-- Sổ tay kiến thức: Bài 42. Số bị chia, số chia, thương --}

\setlength{\parindent}{0pt}
\setlength{\parskip}{0pt}

\begin{document}

% --- BANNER BÀI HỌC ---
\begin{tcolorbox}[
  enhanced,
  colback=violet!8!white,
  colframe=violet!70!black,
  arc=3mm, boxrule=1.2pt,
  left=3mm, right=3mm, top=2mm, bottom=2mm,
  halign=center
]
  {\LARGE\bfseries\color{violet!90!black}\faBookmark\ BÀI 42: SỐ BỊ CHIA, SỐ CHIA, THƯƠNG}\\[1mm]
  {\itshape\footnotesize\color{black!75} Tên gọi các thành phần phép chia và ý nghĩa của thương trên 1 trang A4}
\end{tcolorbox}

\vspace{-2mm}

\noindent
\begin{minipage}[t]{0.49\textwidth}\vspace{0pt}
\begin{tcolorbox}[
  enhanced,
  colback=violet!4!white,
  colframe=violet!75!black,
  arc=2mm, boxrule=0.8pt,
  left=2mm, right=2mm, top=1.8mm, bottom=1.8mm,
  title={\footnotesize\bfseries\color{white}\faKey\ 1. Thành phần phép chia}
]
{\footnotesize
\begin{itemize}\setlength{\itemsep}{1pt}\setlength{\leftskip}{-2.5mm}
  \item \textbf{Cấu trúc tổng quát:}
  $$\mathbf{a : b = c}$$
  \begin{itemize}\setlength{\itemsep}{0.5pt}\setlength{\leftskip}{-2.5mm}
    \item $a$: \textbf{Số bị chia} (số đứng trước dấu chia, đem đi chia).
    \item $b$: \textbf{Số chia} (số đứng sau dấu chia, số phần được chia; $b \neq 0$).
    \item $c$: \textbf{Thương} (kết quả phép chia).
  \end{itemize}
  \item \textbf{Lưu ý đặc biệt:}\\
  Biểu thức $a : b$ cũng được gọi là một \textbf{"thương"}.\\
  \textit{Ví dụ:} Thương của $12$ và $2$ là $12 : 2$ hay $6$.
\end{itemize}
}
\end{tcolorbox}
\end{minipage}
\hfill
\begin{minipage}[t]{0.49\textwidth}\vspace{0pt}
\begin{tcolorbox}[
  enhanced,
  colback=teal!4!white,
  colframe=teal!75!black,
  arc=2mm, boxrule=0.8pt,
  left=2mm, right=2mm, top=1.8mm, bottom=1.8mm,
  title={\footnotesize\bfseries\color{white}\faSliders\ 2. Quan hệ với phép nhân}
]
{\footnotesize
\begin{itemize}\setlength{\itemsep}{1pt}\setlength{\leftskip}{-2.5mm}
  \item Mối liên hệ cốt lõi:
  $$\mathbf{\text{Số bị chia} = \text{Thương} \times \text{Số chia}}$$
  $$\mathbf{\text{Số chia} = \text{Số bị chia} : \text{Thương}}$$
  \item \textbf{Ví dụ:} $10 : 2 = 5$
  \begin{itemize}\setlength{\itemsep}{0.5pt}\setlength{\leftskip}{-2.5mm}
    \item Số bị chia: $10 = 5 \times 2$.
    \item Số chia: $2 = 10 : 5$.
    \item Thương: $5 = 10 : 2$.
  \end{itemize}
\end{itemize}
}
\end{tcolorbox}
\end{minipage}

\vspace{-1.5mm}
\begin{tcolorbox}[
  enhanced,
  colback=orange!4!white,
  colframe=orange!80!black,
  arc=2mm, boxrule=0.8pt,
  left=2mm, right=2mm, top=1.5mm, bottom=1.5mm,
  title={\footnotesize\bfseries\color{white}\faSitemap\ 3. Sơ đồ tư duy định danh các thành phần phép chia}
]
\centering
\begin{tikzpicture}[x=1.5cm, y=0.8cm]
  \node[rectangle, rounded corners=2pt, draw=violet!80!black, fill=violet!15, font=\scriptsize\bfseries] (sbc) at (0,0) {SỐ BỊ CHIA\\($18$)};
  \node[font=\normalsize\bfseries] at (1,0) {$:$};
  \node[rectangle, rounded corners=2pt, draw=teal!80!black, fill=teal!15, font=\scriptsize\bfseries] (sc) at (2,0) {SỐ CHIA\\($2$)};
  \node[font=\normalsize\bfseries] at (3,0) {$=$};
  \node[rectangle, rounded corners=2pt, draw=green!80!black, fill=green!15, font=\scriptsize\bfseries] (th) at (4.5,0) {THƯƠNG\\($9$)};

  \draw[arr, thick, orange!80!black, ->] (th) to[bend left=30] node[above, font=\tiny\bfseries] {Thương $\times$ Số chia $= 9 \times 2 = 18$} (sbc);
\end{tikzpicture}
\end{tcolorbox}

\vspace{-1.5mm}
\noindent
\begin{minipage}[t]{0.49\textwidth}\vspace{0pt}
\begin{tcolorbox}[
  enhanced,
  colback=green!4!white,
  colframe=green!60!black,
  arc=2mm, boxrule=0.8pt,
  left=2mm, right=2mm, top=1.8mm, bottom=1.8mm,
  title={\footnotesize\bfseries\color{white}\faCheckCircle\ 4. Ví dụ thực hành chuẩn mực}
]
{\footnotesize
\textbf{1. Xác định thành phần:}\\
Trong phép tính $16 : 2 = 8$:\\
- $16$ là \textbf{số bị chia}.\\
- $2$ là \textbf{số chia}.\\
- $8$ là \textbf{thương}.
\textbf{2. Tính thương:}\\
Tìm thương của $20$ và $5$?\\
\textit{Giải:} Thương là: $20 : 5 = \mathbf{4}$.
}
\end{tcolorbox}
\end{minipage}
\hfill
\begin{minipage}[t]{0.49\textwidth}\vspace{0pt}
\begin{tcolorbox}[
  enhanced,
  colback=red!4!white,
  colframe=red!75!black,
  arc=2mm, boxrule=0.8pt,
  left=2mm, right=2mm, top=1.8mm, bottom=1.8mm,
  title={\footnotesize\bfseries\color{white}\faExclamationTriangle\ 5. Mẹo tránh bẫy trắc nghiệm}
]
{\footnotesize
\begin{itemize}\setlength{\itemsep}{1pt}\setlength{\leftskip}{-2.5mm}
  \item \textbf{Bẫy từ khóa: "Tìm thương"}\\
  Đề bài: "Tìm thương của 14 và 2" $\implies$ Làm phép \textbf{chia} ($14 : 2 = 7$). Không lấy cộng hay trừ!
  \item Số bị chia đứng \textbf{trước} dấu chia, Số chia đứng \textbf{sau} dấu chia.
  \item Số chia không bao giờ bằng $0$.
\end{itemize}
}
\end{tcolorbox}
\end{minipage}


\end{document}
